\documentclass[10pt,a4paper]{amsart}
\usepackage{amsmath,amssymb,amsthm,graphicx}
% Retained minimal preamble for the Stacks Project chapter sites-cohomology.tex (Cohomology on Sites).
% Verbatim excerpt of the chapter preamble: Xy-pic setup (native preamble lines 14-18) and the native
% theorem declarations with the mathematical macros (native preamble lines 154-217). Copied mechanically;
% nothing added, removed or redefined. The wrapper class and the amsmath/amssymb/amsthm/graphicx base
% packages are supplied by the assembler, so the class/externaldocument part of the native preamble is omitted.
% For commutative diagrams we use Xy-pic
\usepackage[all]{xy}

% We use 2cell for 2-commutative diagrams.
\xyoption{2cell}
% Theorem environments.
%
\theoremstyle{plain}
\newtheorem{theorem}[subsection]{Theorem}
\newtheorem{proposition}[subsection]{Proposition}
\newtheorem{lemma}[subsection]{Lemma}

\theoremstyle{definition}
\newtheorem{definition}[subsection]{Definition}
\newtheorem{example}[subsection]{Example}
\newtheorem{exercise}[subsection]{Exercise}
\newtheorem{situation}[subsection]{Situation}

\theoremstyle{remark}
\newtheorem{remark}[subsection]{Remark}
\newtheorem{remarks}[subsection]{Remarks}

\numberwithin{equation}{subsection}

% Macros
%
\def\lim{\mathop{\mathrm{lim}}\nolimits}
\def\colim{\mathop{\mathrm{colim}}\nolimits}
\def\Spec{\mathop{\mathrm{Spec}}}
\def\Hom{\mathop{\mathrm{Hom}}\nolimits}
\def\Ext{\mathop{\mathrm{Ext}}\nolimits}
\def\SheafHom{\mathop{\mathcal{H}\!\mathit{om}}\nolimits}
\def\SheafExt{\mathop{\mathcal{E}\!\mathit{xt}}\nolimits}
\def\Sch{\mathit{Sch}}
\def\Mor{\mathop{\mathrm{Mor}}\nolimits}
\def\Ob{\mathop{\mathrm{Ob}}\nolimits}
\def\Sh{\mathop{\mathit{Sh}}\nolimits}
\def\NL{\mathop{N\!L}\nolimits}
\def\CH{\mathop{\mathrm{CH}}\nolimits}
\def\proetale{{pro\text{-}\acute{e}tale}}
\def\etale{{\acute{e}tale}}
\def\QCoh{\mathit{QCoh}}
\def\Ker{\mathop{\mathrm{Ker}}}
\def\Im{\mathop{\mathrm{Im}}}
\def\Coker{\mathop{\mathrm{Coker}}}
\def\Coim{\mathop{\mathrm{Coim}}}

% Boxtimes
%
\DeclareMathSymbol{\boxtimes}{\mathbin}{AMSa}{"02}

%
% Macros for moduli stacks/spaces
%
\def\QCohstack{\mathcal{QC}\!\mathit{oh}}
\def\Cohstack{\mathcal{C}\!\mathit{oh}}
\def\Spacesstack{\mathcal{S}\!\mathit{paces}}
\def\Quotfunctor{\mathrm{Quot}}
\def\Hilbfunctor{\mathrm{Hilb}}
\def\Curvesstack{\mathcal{C}\!\mathit{urves}}
\def\Polarizedstack{\mathcal{P}\!\mathit{olarized}}
\def\Complexesstack{\mathcal{C}\!\mathit{omplexes}}
% \Pic is the operator that assigns to X its picard group, usage \Pic(X)
% \Picardstack_{X/B} denotes the Picard stack of X over B
% \Picardfunctor_{X/B} denotes the Picard functor of X over B
\def\Pic{\mathop{\mathrm{Pic}}\nolimits}
\def\Picardstack{\mathcal{P}\!\mathit{ic}}
\def\Picardfunctor{\mathrm{Pic}}
\def\Deformationcategory{\mathcal{D}\!\mathit{ef}}
\usepackage{cleveref}
\crefname{originallemma}{original lemma}{original lemmas}
\Crefname{originallemma}{Original lemma}{Original lemmas}
\makeatletter
\newlabel{sites-section-sheafification}{{00W1}{source}{}{sourcepointer-1}{}}
\newlabel{sites-section-sheafification@cref}{{[originallemma][0][]00W1}{[1][0][]source}{}{}{}}
\makeatother
\crefname{originallemma}{original lemma}{original lemmas}
\Crefname{originallemma}{Original lemma}{Original lemmas}
\makeatletter
\newlabel{sheaves-lemma-sheaf-subset-stalks}{{0079}{source}{}{sourcepointer-2}{}}
\newlabel{sheaves-lemma-sheaf-subset-stalks@cref}{{[originallemma][0][]0079}{[1][0][]source}{}{}{}}
\makeatother
\crefname{originallemma}{original lemma}{original lemmas}
\Crefname{originallemma}{Original lemma}{Original lemmas}
\makeatletter
\newlabel{sheaves-lemma-stalk-pullback-presheaf}{{008G}{source}{}{sourcepointer-3}{}}
\newlabel{sheaves-lemma-stalk-pullback-presheaf@cref}{{[originallemma][0][]008G}{[1][0][]source}{}{}{}}
\makeatother
\crefname{originallemma}{original lemma}{original lemmas}
\Crefname{originallemma}{Original lemma}{Original lemmas}
\makeatletter
\newlabel{sites-theorem-plus}{{00WB}{source}{}{sourcepointer-4}{}}
\newlabel{sites-theorem-plus@cref}{{[originallemma][0][]00WB}{[1][0][]source}{}{}{}}
\makeatother
\crefname{originallemma}{original lemma}{original lemmas}
\Crefname{originallemma}{Original lemma}{Original lemmas}
\makeatletter
\newlabel{sheaves-section-sheafification}{{007X}{source}{}{sourcepointer-5}{}}
\newlabel{sheaves-section-sheafification@cref}{{[originallemma][0][]007X}{[1][0][]source}{}{}{}}
\makeatother
\title{The pullback of a sheaf along the projection of the qc site of locally quasi-compact spaces}
\author{The Stacks Project Authors}
\date{stacks-project:sites-cohomology.tex:local-snapshot-20261011; GNU FDL 1.2 or later}
\begin{document}
\maketitle
\noindent\textit{Editorial note.} All delivered mathematical text is the author's own, copied from the pinned Stacks Project chapter file sites-cohomology.tex; no statement or proof step was rewritten, reordered or removed.\par\medskip
\noindent\textit{Editorial note.} The proof is the author's complete argument for both assertions of the lemma: the sheaf property in the qc topology and the identification with the two pullbacks are each carried through inside the counted span.\par\medskip
\noindent\textit{Editorial note.} Every labelled reference occurring inside the counted spans points to a chapter of the Project that is not part of this card's source tree (Sheaves and Sites); each is carried in assembly.external\_references as a non-load-bearing navigation pointer whose source location is the invocation inside the counted span, so no load-bearing mathematics is replaced by a stub.\par\medskip
\noindent\textit{Editorial note.} The checked chapter list for same-name collisions is the set of stacks chapters staged in this repository: chow, cohomology, divisors, etale-cohomology, flat, more-algebra, more-morphisms, morphisms, spaces-more-morphisms, varieties together with this card's chapter sites-cohomology.\par\medskip
\noindent\textit{Editorial note.} The preamble is the chapter's own retained minimal preamble (sources/stacks-sites-cohomology/retained-preamble.tex), a verbatim excerpt of the native preamble: the Xy-pic setup and the native theorem declarations with the mathematical macros.\par\medskip
\noindent\textit{Allowed standard prerequisites (editorial).}\begin{itemize}
\item The category LC\_qc of locally quasi-compact spaces and the qc topology. Same chapter, the category LC of Hausdorff locally quasi-compact spaces with the qc topology, and the morphisms pi\_X and epsilon\_X comparing it with the Zariski topology. Supplies the sites between which the pullbacks in the statement are computed.
\item qc covering and the sheaf condition for the qc topology. Same chapter, the definition of a qc covering and the criterion that a presheaf is a sheaf once its sections over every object map isomorphically to the degree-zero Cech cohomology of every qc covering. Supplies the reduction of the sheaf condition to the injectivity and gluing statements proved in the second part of the proof.
\item Separated presheaves and the plus construction. Sites, the theorem that a separated presheaf becomes a sheaf after applying the plus construction once, together with the material on sheafification and stalks. Supplies the reduction from gluing after refinement to gluing, and the comparison of sections by stalks.
\item Pullback of sheaves on topological spaces and its stalks. Sheaves, the description of f\textasciicircum{}\{-1\}F by the presheaf colim over the images of open sets, and the computation of its stalks. Supplies the formula for the pullback used in the second and third parts of the proof.
\end{itemize}\medskip
\ifdefined\hypertarget\hypertarget{sourcepointer-1}{}\fi
\noindent\textit{Source reference.} Original lemma 00W1: sources/stacks-sites-cohomology/pinned-original-source.tex, source lines 7701--7701. Navigation pointer invoked at proof line 7701; the counted span contains the pointer only and the cited result is not reproduced inside it, so the pointer is registered as a non-load-bearing navigation stub in the established form for Stacks material. No bibliography citation is involved.\par\medskip
\ifdefined\hypertarget\hypertarget{sourcepointer-2}{}\fi
\noindent\textit{Source reference.} Original lemma 0079: sources/stacks-sites-cohomology/pinned-original-source.tex, source lines 7704--7704. Navigation pointer invoked at proof line 7704; the counted span contains the pointer only and the cited result is not reproduced inside it, so the pointer is registered as a non-load-bearing navigation stub in the established form for Stacks material. No bibliography citation is involved.\par\medskip
\ifdefined\hypertarget\hypertarget{sourcepointer-3}{}\fi
\noindent\textit{Source reference.} Original lemma 008G: sources/stacks-sites-cohomology/pinned-original-source.tex, source lines 7705--7705. Navigation pointer invoked at proof line 7705; the counted span contains the pointer only and the cited result is not reproduced inside it, so the pointer is registered as a non-load-bearing navigation stub in the established form for Stacks material. No bibliography citation is involved.\par\medskip
\ifdefined\hypertarget\hypertarget{sourcepointer-4}{}\fi
\noindent\textit{Source reference.} Original lemma 00WB: sources/stacks-sites-cohomology/pinned-original-source.tex, source lines 7711--7711. Navigation pointer invoked at proof line 7711; the counted span contains the pointer only and the cited result is not reproduced inside it, so the pointer is registered as a non-load-bearing navigation stub in the established form for Stacks material. No bibliography citation is involved.\par\medskip
\ifdefined\hypertarget\hypertarget{sourcepointer-5}{}\fi
\noindent\textit{Source reference.} Original lemma 007X: sources/stacks-sites-cohomology/pinned-original-source.tex, source lines 7736--7736. Navigation pointer invoked at proof line 7736; the counted span contains the pointer only and the cited result is not reproduced inside it, so the pointer is registered as a non-load-bearing navigation stub in the established form for Stacks material. No bibliography citation is involved.\par\medskip
\begin{lemma}
\label{lemma-describe-pullback-pi}
Let $X$ be an object of $\textit{LC}_{qc}$. Let $\mathcal{F}$ be a
sheaf on $X$. The rule
$$
\textit{LC}_{qc}/X \longrightarrow \textit{Sets},\quad
(f : Y \to X) \longmapsto \Gamma(Y, f^{-1}\mathcal{F})
$$
is a sheaf and a fortiori also a sheaf on $\textit{LC}_{Zar}/X$.
This sheaf is equal to
$\pi_X^{-1}\mathcal{F}$ on $\textit{LC}_{Zar}/X$ and
$\epsilon_X^{-1}\pi_X^{-1}\mathcal{F}$ on $\textit{LC}_{qc}/X$.
\end{lemma}

\begin{proof}
Denote $\mathcal{G}$ the presheaf given by the formula in the lemma.
Of course the pullback $f^{-1}$ in the formula denotes usual
pullback of sheaves on topological spaces. It is immediate
from the definitions that $\mathcal{G}$ is a sheaf for the Zar
topology.

\medskip\noindent
Let $Y \to X$ be a morphism in $\textit{LC}_{qc}$. Let
$\mathcal{V} = \{g_i : Y_i \to Y\}_{i \in I}$ be a qc covering.
To prove $\mathcal{G}$ is a sheaf for the qc topology it
suffices to show that
$\mathcal{G}(Y) \to H^0(\mathcal{V}, \mathcal{G})$
is an isomorphism, see Sites, Section \ref{sites-section-sheafification}.
We first point out that the map is injective as a qc covering
is surjective and we can detect equality of sections at stalks
(use Sheaves, Lemmas \ref{sheaves-lemma-sheaf-subset-stalks} and
\ref{sheaves-lemma-stalk-pullback-presheaf}). Thus
$\mathcal{G}$ is a separated presheaf on $\textit{LC}_{qc}$
hence it suffices to show that any element
$(s_i) \in H^0(\mathcal{V}, \mathcal{G})$
maps to an element in the image of $\mathcal{G}(Y)$
after replacing $\mathcal{V}$ by a refinement
(Sites, Theorem \ref{sites-theorem-plus}).

\medskip\noindent
Identifying sheaves on $Y_{i, Zar}$ and sheaves on $Y_i$ we find that
$\mathcal{G}|_{Y_{i, Zar}}$ is the pullback of $f^{-1}\mathcal{F}$ under
the continuous map $g_i : Y_i \to Y$. Thus we can choose an open covering
$Y_i = \bigcup V_{ij}$ such that for each $j$ there is an open
$W_{ij} \subset Y$ and a section $t_{ij} \in \mathcal{G}(W_{ij})$
such that $V_{ij}$ maps into $W_{ij}$ and such that
$s|_{V_{ij}}$ is the pullback of $t_{ij}$. In other words,
after refining the covering $\{Y_i \to Y\}$ we may assume there
are opens $W_i \subset Y$ such that $Y_i \to Y$ factors through $W_i$
and sections $t_i$ of $\mathcal{G}$ over $W_i$ which restrict
to the given sections $s_i$. Moreover, if $y \in Y$ is in the image
of both $Y_i \to Y$ and $Y_j \to Y$, then the images $t_{i, y}$
and $t_{j, y}$ in the stalk $f^{-1}\mathcal{F}_y$ agree
(because $s_i$ and $s_j$ agree over $Y_i \times_Y Y_j$).
Thus for $y \in Y$ there is a well defined element $t_y$ of
$f^{-1}\mathcal{F}_y$ agreeing with $t_{i, y}$ whenever $y$
is in the image of $Y_i \to Y$.
We will show that the element $(t_y)$ comes from a global section
of $f^{-1}\mathcal{F}$ over $Y$ which will finish the proof of the lemma.

\medskip\noindent
It suffices to show that this is true locally on $Y$, see
Sheaves, Section \ref{sheaves-section-sheafification}. Let $y_0 \in Y$.
Pick $i_1, \ldots, i_n \in I$ and
quasi-compact subsets $E_j \subset Y_{i_j}$ such that
$\bigcup g_{i_j}(E_j)$ is a neighbourhood of $y_0$.
Let $V \subset Y$ be an open neighbourhood of $y_0$ contained
in $\bigcup g_{i_j}(E_j)$ and contained in $W_{i_1} \cap \ldots \cap W_{i_n}$.
Since $t_{i_1, y_0} = \ldots = t_{i_n, y_0}$, after shrinking $V$
we may assume the sections $t_{i_j}|_V$, $j = 1, \ldots, n$ of
$f^{-1}\mathcal{F}$ agree. As $V \subset \bigcup g_{i_j}(E_j)$
we see that $(t_y)_{y \in V}$ comes from this section.

\medskip\noindent
We still have to show that $\mathcal{G}$ is equal to
$\epsilon_X^{-1}\pi_X^{-1}\mathcal{F}$ on $\textit{LC}_{qc}$,
resp.\ $\pi_X^{-1}\mathcal{F}$ on $\textit{LC}_{Zar}$.
In both cases the pullback is defined by taking the presheaf
$$
(f : Y \to X)
\longmapsto
\colim_{f(Y) \subset U \subset X} \mathcal{F}(U)
$$
and then sheafifying. Sheafifying in the Zar topology
exactly produces our sheaf $\mathcal{G}$ and the fact
that $\mathcal{G}$ is a qc sheaf, shows that it works as well
in the qc topology.
\end{proof}

\end{document}
